\documentclass[10pt]{article}
\usepackage[T1]{fontenc}
\usepackage[utf8]{inputenc}
\usepackage[french]{babel}
\usepackage{amsmath,amssymb,mathrsfs,mathpazo}
\usepackage{geometry}
\geometry{a4paper,margin=2cm}
\usepackage{graphicx,array,tabularx,tabularray,colortbl,multirow,multicol,booktabs,enumitem,pifont,ulem,fancybox,diagbox,cancel}
\usepackage{pstricks,pstricks-add,pst-plot,pst-tree,pst-node,pst-text,pst-func,pst-math,pst-bspline,pst-3dplot,pst-eucl,pst-solides3d}
\let\pscalculate\undefined
\usepackage{pst-bezier}
\usepackage{tikz,pgfplots,tkz-tab}
\usetikzlibrary{arrows,arrows.meta,patterns,calc,positioning,trees,decorations.markings,shapes.geometric,shapes.misc,shapes.arrows,intersections,decorations.pathmorphing,decorations.pathreplacing,backgrounds,shadows,babel}
\pgfplotsset{compat=1.18}
\usepgfplotslibrary{fillbetween}
\usepackage[np]{numprint}
\usepackage{listings,makecell,pas-tableur,twoopt}
\let\barre\undefined
\usepackage[pstricks]{tablvar}
\definecolor{darkgreen}{RGB}{0,100,0}
\providecommand{\red}{\color{red}}
\providecommand{\blue}{\color{blue}}
\providecommand{\green}{\color{green}}
\providecommand{\black}{\color{black}}
\providecommand{\white}{\color{white}}
\providecommand{\gray}{\color{gray}}
\providecommand{\cyan}{\color{cyan}}
\providecommand{\magenta}{\color{magenta}}
\providecommand{\yellow}{\color{yellow}}
\providecommand{\violet}{\color{violet}}
\providecommand{\orange}{\color{orange}}
\providecommand{\euro}{\texteuro}
\providecommand{\EUR}{\texteuro}
\providecommand{\ieme}{\textsuperscript{e}}
\providecommand{\degres}{\ensuremath{^\circ}}
\providecommand{\degree}{\ensuremath{^\circ}}
\providecommand{\vv}[1]{\overrightarrow{#1}}
\providecommand{\Coord}[1]{#1}
\providecommand{\decosix}{}
\providecommand{\starredbullet}{\ensuremath{\bullet}}
\pagestyle{empty}
\setlength{\parindent}{0pt}
\usepackage{adjustbox}\DeclareUnicodeCharacter{2236}{:}\DeclareUnicodeCharacter{2190}{\ensuremath{\leftarrow}}
\providecommand{\R}{\mathbb{R}}
\providecommand{\N}{\mathbb{N}}
\providecommand{\Z}{\mathbb{Z}}
\providecommand{\Q}{\mathbb{Q}}
\providecommand{\D}{\mathbb{D}}
\providecommand{\C}{\mathbb{C}}
\providecommand{\e}{{\mathrm{e}}}
\providecommand{\vect}[1]{\overrightarrow{#1}}
\providecommand{\vectt}[1]{\overrightarrow{\mathrm{#1}}}
\providecommand{\vvt}[1]{\overrightarrow{\mathrm{#1}}}
\providecommand{\np}[2][]{{#2~#1}}
\providecommand{\vv}[1]{{\overrightarrow{#1}}}
\renewcommand{\d}{{\mathrm{d}}}
\providecommand{\degres}{{^\circ}}
\providecommand{\degre}{{^\circ}}
\providecommand{\blue}{\color{blue}}
\providecommand{\red}{\color{red}}
\providecommand{\enskip}{\quad}
\providecommand{\llbracket}{\lbrack\!\lbrack}
\providecommand{\rrbracket}{\rbrack\!\rbrack}
\providecommand{\euro}{{\text{€}}}
\providecommand{\no}{{\text{n°}}}
\providecommand{\interval}[2]{{[#1;#2]}}
\providecommand{\intervFO}[2]{{[#1;#2[}}
\providecommand{\dotfill}{\ldots}

\providecommand{\ds}{\displaystyle}
\providecommand{\textding}[1]{#1}
\providecommand{\Oij}{\(\mathrm{O};\vec{\imath},\vec{\jmath}\)}
\providecommand{\Oijk}{\(\mathrm{O};\vec{\imath},\vec{\jmath},\vec{k}\)}
\providecommand{\Ouv}{\(\mathrm{O};\vec{u},\vec{v}\)}
\providecommand{\intervFF}[2]{\ensuremath{[#1;#2]}}
\renewcommand{\euro}{\ifmmode\text{\texteuro}\else\texteuro\fi}
\DeclareUnicodeCharacter{00B0}{\ifmmode^{\circ}\else\textdegree\fi}
\DeclareUnicodeCharacter{20AC}{\ifmmode\text{\texteuro}\else\texteuro\fi}
% Énoncés : Éducation nationale. Transcription et corrigés : APMEP, Denis Vergès et contributeurs. Découpage et classement : CoopMaths, Sébastien Lozano, Éric Elter et contributeurs. Adaptation de la présentation : Maths Explorer.
% CC BY-SA : https://coopmaths.fr/www/about/
% https://coopmaths.fr/alea/static/bac/2025/tex/bac_2025_11_sujet2_ameriquesud_1_cor.tex
\begin{document}
\noindent\textit{APMEP / CoopMaths. CC BY-SA. Présentation adaptée par Maths Explorer.}\par\medskip
\medskip

\emph{Dans cet exercice, tous les résultats seront arrondis à $10^{-3}$ près en cas de besoin.\\
Les deux parties de cet exercice sont indépendantes l'une de l'autre.}

\bigskip

\textbf{Partie A}

\medskip

Au tennis, le joueur qui est au service peut, en cas d'échec lors du premier service, servir une deuxième balle.

En match, Abel réussit son premier service dans 70\,\% des cas. Lorsque le premier service est réussi, il gagne le point dans 80\,\% des cas.

En revanche, après un échec à son premier service, Abel gagne
le point dans 45\,\% des cas.



Abel est au service.

\begin{list}{\textbullet}{On considère les évènements suivants:}
\item $S$ : \og Abel réussit son premier service \fg{};
\item $G$ : \og Abel gagne le point \fg.
\end{list}



\begin{enumerate}
\item $\overline{S}$ est l'événement: \og Abel ne réussit pas son premier service. \fg{}.

On traduit la situation par un arbre pondéré.

\begin{center}

  \pstree[treemode=R,nodesepA=0pt,nodesepB=4pt,levelsep=2.5cm,treesep=1cm]{\TR{}}
 {
 	\pstree[nodesepA=4pt]{\TR{$S$}\naput{$0,7$}}
 	  { 
 		  \TR{$G$}\naput{$0,8$}
 		  \TR{$\overline{G}$}\nbput{$0,2$}	   
 	  }
 	\pstree[nodesepA=4pt]{\TR{$\overline{S}$}\nbput{$0,3$}}
 	  {
 		  \TR{$G$}\naput{$0,45$}
 		  \TR{$\overline{G}$}\nbput{$0,55$}	   
      }
}

\end{center}	

\item $P\left (S \cap G\right ) = P\left (S\right ) \times P_{S}\left (G\right ) = 0,7\times 0,8 = 0,56$

\item 
$\left \{ S\,,\, \overline{S}\strut \right \}$ forme une partition de l'univers donc, d'après la formule des probabilités totales, on a:
$P\left (G\right ) = P\left (S\cap G\right ) + P \left (\overline{S}\cap G\right )
= 0,56 + 0,3\times 0,45 = 0,695$

\item Abel a gagné le point. La probabilité qu'il ait réussi son premier service est:

$P_{G}\left (S\right ) = \dfrac{P\left (S\cap G \right )}{P\left (G\right )}=\dfrac{0,56}{0,695}\approx 0,806$

\item 
\begin{list}{\textbullet}{On a:}
\item $P\left (S\cap G\right )=0,56$
\item $P\left (S\right )\times P\left (G\right ) = 0,7 \times 0,695 = \np{0,4865}$
\end{list}

$P\left (S\cap G\right ) \neq  P\left (S\right )\times P\left (G\right )$ donc les événements $S$ et $G$ ne sont pas indépendants.

\end{enumerate}


\textbf{Partie B}

\medskip

À la sortie d'une usine de fabrication de balles de tennis, une balle est jugée conforme dans 85\,\% des cas.



\begin{enumerate}
\item On teste successivement $20$~balles. On considère que le nombre de balles est suffisamment grand pour assimiler ces tests à un tirage avec remise. On note $X$ la variable aléatoire qui compte le nombre de balles conformes parmi les $20$ testées.
	\begin{enumerate}
		\item 
L'épreuve élémentaire consiste à voir si une balle est conforme, avec une probabilité de $0,85$, ou si elle ne l'est pas.

On exécute cette épreuve élémentaire 20 fois et on  considère que le nombre de balles est suffisamment grand pour assimiler ces tests à un tirage avec remise.

Donc la variable aléatoire $X$ qui compte le nombre de balles conformes parmi les $20$ testées, suit la loi binomiale de paramètres $n=20$ et $p=0,85$.
		
		\item À la calculatrice,  $P\left (X \leqslant 18\right ) \approx 0,824$.
		
		\item La probabilité qu'au moins deux balles ne soient pas conformes parmi les $20$ balles testées est  égale à la probabilité qu'au plus 18 balles soient conformes, c'est-à-dire $P\left (X\leqslant18\right )$ soit environ $0,824$.
		
		
		\item L'espérance de $X$ est $E(X)=np=20\times 0,85=17$.
	\end{enumerate}
	
\item On teste maintenant $n$ balles successivement. On considère les $n$ tests comme un échantillon de $n$ variables aléatoires $X$ indépendantes suivant la loi de Bernoulli de paramètre $0,85$.
	
On considère la variable aléatoire
	$M_n = \displaystyle\sum_{i=1}^n \dfrac{X_i}{n} = \dfrac{X_1}{n} + \dfrac{X_2}{n} + \dfrac{X_3}{n} + \ldots + \dfrac{X_n}{n}$

	\begin{enumerate}
		\item 
D'après le cours, on peut dire que les  $n$ variables aléatoires $X$ indépendantes qui suivent la loi de Bernoulli de paramètre $p=0,85$ ont pour espérance $p=0,85$ et pour variance $p\left (1-p\right )=0,85\times 0,15=\np{0,1275}$.

\begin{list}{\textbullet}{Soit $S_n$ la loi somme   définie par 	$S_n=X_1+X_2+X_3+ \cdots + X_n  = \ds\sum_{i=1}^n X_i$.}
\item D'après la linéarité de l'espérance, on a:

$E\left (S_n\right )= E\left (\ds\sum_{i=1}^n X_i \right ) = \ds\sum_{i=1}^n E\left (X_i\right ) = \ds\sum_{i=1}^n 0,85=0,85n$
\item Les variables $X_i$ étant indépendantes, on utilise l'additivité de la variance:

$V\left (S_n\right )= V\left (\ds\sum_{i=1}^n X_i \right ) = \ds\sum_{i=1}^n V\left (X_i\right ) = \ds\sum_{i=1}^n \np{0,1275}=\np{0,1275}n$
\end{list}



\begin{list}{\textbullet}{$M_n$ est la loi moyenne définie par $M_n=\dfrac{S_n}{n}$.}
\item On sait que  $E(aX)=aE(X)$, donc:

$E\left (M_n\right )=E\left (\dfrac{1}{n}S_n\right ) = \dfrac{1}{n} \; E\left (S_n\right ) = \dfrac{0,85n}{n}=0,85$.
\item On sait que $V(aX)=a^2V(X)$ donc:

$V\left (M_n\right )=V\left (\dfrac{1}{n}S_n\right ) = \dfrac{1}{n^2} \; V\left (S_n\right ) = \dfrac{\np{0,1275}n}{n^2}= \dfrac{\np{0,1275}}{n}$.
\end{list}
				
		\item 
Si $X$ est une variable aléatoire et $t$ un réel strictement positif, on a:

$P\left ( \left | X-E(X)\strut \right | \geqslant t\right ) \leqslant \dfrac{V(X)}{t^2}$.
C'est l'inégalité de Bienaymé-Tchebychev.

On en déduit que: 	$P\left ( \left | X-E(X)\strut \right | < t\right ) \geqslant 1- \dfrac{V(X)}{t^2}$. Donc:
		
$\begin{aligned}
P\left ( \left | X-E(X)\strut \right | < t\right ) \geqslant 1- \dfrac{V(X)}{t^2}
& \iff
P\left ( \left | M_n-E(M_n)\strut \right | < t\right ) \geqslant 1- \dfrac{V(M_n)}{t^2}\\
& \iff
P\left ( \left | M_n-0,85\strut \right | < t\right ) \geqslant 1- \dfrac{\np{0,1275}}{nt^2}
\end{aligned}$

\begin{list}{\textbullet}{On prend $t=0,1$. On a donc:}
\item 
$\begin{aligned}[t]
\left | M_n-0,85\strut \right | < t
& \iff \left | M_n-0,85\strut \right | < 0,1
\iff -0,1 <  M_n-0,85< 0,1\\
& \iff 0,75 < M_n < 0,95
\end{aligned}$

\item $1- \dfrac{\np{0,1275}}{nt^2} =  1- \dfrac{\np{0,1275}}{n\times 0,1^2} =  1- \dfrac{12,75}{n}$
\end{list} 

On déduit donc:
$P(0,75 < M_n < 0,95) \geqslant  1 - \dfrac{12,75}{n}$.

		\item Pour trouver un entier $n$ tel que la moyenne du nombre de balles conformes pour un échantillon de taille $n$ appartienne à l'intervalle ]0,75~;~0,95 [ avec une probabilité supérieure à $0,9$, il suffit que: $1 - \dfrac{12,75}{n}\geqslant 0,9$.
On résout cette inéquation.

$1 - \dfrac{12,75}{n}\geqslant 0,9
 \iff 0,1\geqslant \dfrac{12,75}{n}
\iff 0,1n \geqslant 12,75
\iff n \geqslant 127,5$

Donc il faut un échantillon de taille supérieure à $n=128$ pour  que  la moyenne du nombre de balles conformes  appartenant à l'intervalle $]0,75~;~0,95[$ ait une probabilité supérieure à $0,9$
	\end{enumerate}
\end{enumerate}

\vspace{0,5cm}

\end{document}
